% Chapter 14: Tense and Time

\chapter{Tense and Time: When Does It Happen?}

\section{Pedagogical Purpose}
Having mastered subject-verb agreement, learners are ready to explore how verbs also change to show *when* something happens. Tense is one of the most important tools for clear communication, and errors in tense are among the most common mistakes young writers make. This chapter builds a solid foundation in the three main time zones of English before later chapters refine how tenses are used in connected writing.

\begin{learningobjectives}
The learner can:
\begin{itemize}
    \item identify whether a sentence is in the present, past, or future tense.
    \item explain how regular verbs change to show past time.
    \item use irregular past tense forms of common verbs correctly.
    \item distinguish tense by matching it to time expressions.
    \item correct sentences with the wrong tense.
\end{itemize}
\end{learningobjectives}

\section{Prerequisite Check}
\begin{quickcheck}
Circle the verb in each sentence.
\begin{enumerate}
    \item Maya \textcircled{plays} the piano every day.
    \item Rohan \textcircled{finished} his homework.
\end{enumerate}
\end{quickcheck}

\begin{rememberbox}
A verb tells us about an action or a state (Chapter 8). Now we will learn that verbs also carry a clue about \textbf{time}.
\end{rememberbox}

\section{Chapter Opener}
Leo the parrot has been given a diary by Ms. Sharma, and he is very confused about what to write.

\textbf{Leo:} "Squawk! Yesterday I \textbf{fly} to the window! Tomorrow I \textbf{flew} to the window! Today I \textbf{will fly}— no wait, that's not right either!"

\textbf{Maya:} (laughing) "Leo, you've mixed up your times! You need different verb forms for yesterday, today, and tomorrow."

\textbf{Rohan:} "That's tense, isn't it? Ms. Sharma told us about that!"

\textbf{Ms. Sharma:} "Exactly, Rohan! Let's help Leo write his diary correctly by learning about \textbf{tense} — how verbs show time."

\section{Language Experience}
\textbf{Leo's Muddled Diary}

> Yesterday: I fly to the window and I eat a seed.
> Today: I sing a song and I sat on my perch.
> Tomorrow: I flew around the room.

Something is wrong in every line!

\section{Look and Notice}
\begin{thinkbox}
    \begin{enumerate}
        \item Which time word tells you Leo is talking about the past?
        \item Which time word tells you Leo is talking about the future?
        \item Does "fly" sound right for something that already happened?
        \item Can you spot a verb that is in the wrong time zone?
    \end{enumerate}
\end{thinkbox}

\section{Core Explanation}
Verbs change their form to show \textbf{when} an action happens. This is called \textbf{tense}. There are three main time zones:

\begin{table}[h!]
\centering
\begin{tabular}{|l|l|l|}
\hline
\textbf{Time Zone} & \textbf{Meaning} & \textbf{Example} \\
\hline
Present & Happening now, or always true, or a habit & Maya \textbf{plays} the piano. \\
Past & Already happened & Maya \textbf{played} the piano yesterday. \\
Future & Will happen later & Maya \textbf{will play} the piano tomorrow. \\
\hline
\end{tabular}
\end{table}

\subsection*{Regular Verbs (Past Tense)}
For most verbs, we form the past tense by adding \textbf{-ed}.
\begin{itemize}
    \item walk \(\rightarrow\) walk\textbf{ed}
    \item play \(\rightarrow\) play\textbf{ed}
    \item look \(\rightarrow\) look\textbf{ed}
\end{itemize}

\subsection*{Irregular Verbs (Past Tense)}
Some verbs do not follow the -ed rule. They change their form completely.
\begin{itemize}
    \item go \(\rightarrow\) \textbf{went}
    \item eat \(\rightarrow\) \textbf{ate}
    \item see \(\rightarrow\) \textbf{saw}
    \item is/am/are \(\rightarrow\) \textbf{was/were}
\end{itemize}

\begin{warningbox}
    Learners often add \textbf{-ed} to irregular verbs, creating words like "goed" or "eated". This is incorrect. Irregular verbs must be memorized.
\end{warningbox}

\section{Worked Example}
\begin{examplebox}
\textbf{Sentence:} "Yesterday, Maya ___ (go) to the library and ___ (read) two books."

\textbf{Let's Think Step-by-Step:}
\begin{enumerate}
    \item Look for a time clue. "Yesterday" tells us this is the past.
    \item Check each verb. "Go" is irregular in the past — it becomes "went," not "goed."
    \item Check the second verb. "Read" is regular in spelling but pronounced differently; its past form is also "read" (spelled the same, said differently).
\end{enumerate}

\textbf{Answer:} "Yesterday, Maya \textbf{went} to the library and \textbf{read} two books."

\textbf{Why:} The time word "yesterday" signals past tense, so both verbs must match that time zone.
\end{examplebox}
\section{Guided Practice}
\begin{activitybox}{Activity A: Identify the Tense}
Underline the verb and state its tense (Present, Past, or Future).
\begin{enumerate}
    \item Maya \underline{sings} in the choir. (Present)
    \item Rohan \underline{jumped} over the puddle. (Past)
    \item We \underline{will travel} to Delhi next month. (Future)
    \item I \underline{ate} an apple. (Past)
    \item They \underline{play} outside. (Present)
\end{enumerate}
\end{activitybox}

\begin{activitybox}{Activity B: Completion}
Fill in the correct past tense form.
\begin{enumerate}
    \item She (eat) \_\_\_ breakfast at seven.
    \item They (see) \_\_\_ a rainbow after the rain.
\end{enumerate}
\end{activitybox}

\section{Independent Practice}
\begin{activitybox}{Activity B: Change the Tense}
Rewrite the sentences in the tense indicated.
\begin{enumerate}
    \item I write a story. (Change to Past Tense)
    \item Leo eats seeds. (Change to Future Tense)
    \item They play in the park. (Change to Past Tense)
    \item She will sing a song. (Change to Present Tense)
\end{enumerate}
\end{activitybox}

\section{Summary}
\begin{summarybox}
    \begin{itemize}
        \item \textbf{Tense} shows \textbf{when} an action happens: present, past, or future.
        \item \textbf{Regular verbs} add \textbf{-ed} for the past tense (walk \(\rightarrow\) walked).
        \item \textbf{Irregular verbs} change completely in the past tense and must be memorised (go \(\rightarrow\) went).
        \item Time words like \textit{yesterday}, \textit{now}, and \textit{tomorrow} are strong clues for choosing the correct tense.
    \end{itemize}
\end{summarybox}